\section{Alien Dictionary}
\label{sec:graph-alien-dictionary}

\problemstatement{Given a list of words sorted lexicographically
according to the rules of an unknown alien language with $k$
distinct lowercase letters, determine a valid ordering of the
alphabet. Return an empty string if no valid ordering exists.}

\begin{patternbox}
This closing problem of the chapter is a two-step composition:
\textbf{build} a directed graph of letter-ordering constraints by
comparing each pair of \textbf{adjacent} words in the list, then
\textbf{run topological sort} (Sections~\ref{sec:graph-topo-sort-dfs}--\ref{sec:graph-topo-sort-kahn})
on that graph to recover the alphabet order. The graph-building step is
the only new idea; the ordering step is unchanged from earlier in this
chapter.
\end{patternbox}

\subsection*{Extracting constraints from adjacent word pairs}

For two adjacent words $w_1, w_2$ in the sorted list, scan character by
character until the \textbf{first} position where they differ. If
$w_1[i] \neq w_2[i]$, that single comparison reveals exactly one
constraint: $w_1[i]$ comes before $w_2[i]$ in the alien alphabet
(edge $w_1[i] \to w_2[i]$) -- and nothing more can be concluded from
this pair, since the relative order of any later characters is
unconstrained by this comparison. If no differing position is found
within the shorter word's length, the ordering is \textbf{invalid}
exactly when $w_1$ is \emph{longer} than $w_2$ (a longer word cannot
correctly precede its own prefix in sorted order).

\subsection*{Algorithm}

\begin{enumerate}
  \item Initialize an empty adjacency structure over the (at most
        $26$) distinct letters appearing in the words.
  \item For every adjacent pair of words, find the first differing
        character and add the corresponding edge; if no differing
        character exists and the first word is longer, return
        \texttt{""} immediately (invalid).
  \item Run topological sort (Kahn's algorithm is convenient here,
        since it naturally detects a cycle -- contradictory
        constraints -- by an incomplete processed count) over the
        graph of letters.
  \item If a cycle is detected, return \texttt{""}; otherwise return
        the topological order as a string.
\end{enumerate}

\subsection*{C++ implementation}

\begin{lstlisting}
#include <bits/stdc++.h>
using namespace std;

string alienOrder(vector<string>& words) {
    unordered_set<char> letters;
    for (string& w : words) for (char c : w) letters.insert(c);

    unordered_map<char, unordered_set<char>> adj;
    unordered_map<char, int> inDegree;
    for (char c : letters) inDegree[c] = 0;

    for (int i = 0; i + 1 < (int)words.size(); i++) {
        string& w1 = words[i];
        string& w2 = words[i + 1];
        int minLen = min(w1.size(), w2.size());
        bool foundDiff = false;

        for (int j = 0; j < minLen; j++) {
            if (w1[j] != w2[j]) {
                if (adj[w1[j]].find(w2[j]) == adj[w1[j]].end()) {
                    adj[w1[j]].insert(w2[j]);
                    inDegree[w2[j]]++;
                }
                foundDiff = true;
                break;
            }
        }
        if (!foundDiff && w1.size() > w2.size()) return "";
    }

    queue<char> q;
    for (char c : letters) {
        if (inDegree[c] == 0) q.push(c);
    }

    string order;
    while (!q.empty()) {
        char c = q.front();
        q.pop();
        order += c;

        for (char neighbor : adj[c]) {
            inDegree[neighbor]--;
            if (inDegree[neighbor] == 0) q.push(neighbor);
        }
    }

    return (int)order.size() == (int)letters.size() ? order : "";
}

int main() {
    vector<string> words = {"wrt", "wrf", "er", "ett", "rftt"};
    cout << alienOrder(words) << endl; // wertf
    return 0;
}
\end{lstlisting}

\subsection*{Dry run}

\dryrun{Take $\textit{words}=[\texttt{wrt},\texttt{wrf},\texttt{er},
\texttt{ett},\texttt{rftt}]$.}

Pair (\texttt{wrt},\texttt{wrf}): differ at index $2$ ($\texttt{t}$
vs $\texttt{f}$) -- edge $t \to f$. Pair (\texttt{wrf},\texttt{er}):
differ at index $0$ ($\texttt{w}$ vs $\texttt{e}$) -- edge $w \to e$.
Pair (\texttt{er},\texttt{ett}): differ at index $1$ ($\texttt{r}$ vs
$\texttt{t}$) -- edge $r \to t$. Pair (\texttt{ett},\texttt{rftt}):
differ at index $0$ ($\texttt{e}$ vs $\texttt{r}$) -- edge $e \to r$.
Graph: $t\to f$, $w\to e$, $r\to t$, $e\to r$. In-degrees: $w$:$0$,
$e$:$1$, $r$:$1$, $t$:$1$, $f$:$1$. Kahn's: start $[w]$. Pop $w$:
order $\texttt{w}$; $e\to0$, push. Pop $e$: order $\texttt{we}$;
$r\to0$, push. Pop $r$: order $\texttt{wer}$; $t\to0$, push. Pop $t$:
order $\texttt{wert}$; $f\to0$, push. Pop $f$: order
$\texttt{wertf}$. All $5$ letters processed: valid order
$\texttt{"wertf"}$.

\begin{complexitybox}
\textbf{Time Complexity.} $O(\sum |w_i| + V + E)$ -- comparing each
adjacent word pair costs time proportional to the shorter word's
length, and the topological sort afterward costs $O(V+E)$ over the
(at most $26$) letters and their constraints.
\textbf{Space Complexity.} $O(V+E)$ for the letter graph, bounded by
$O(26^2)$ in the worst case.
\end{complexitybox}

\begin{takeawaybox}
This chapter closes by demonstrating topological sort's full
generality: the ``nodes'' and ``edges'' don't need to be given
explicitly at all -- here they are \emph{derived} from a completely
different kind of input (a sorted word list) via a simple
pairwise-comparison rule, and once that translation step is done, every
algorithm from earlier in this chapter (cycle detection bundled into
Kahn's algorithm, the topological order itself) applies completely
unchanged.
\end{takeawaybox}
